\section{Scope, limitations, and future work}
\label{sec:scope}

\subsection{What is established}
The checked endpoint is concrete: a constructive complex commutative ring
$\CH$, an algebraic group of units $\CHunits$, its regular left torsor, the
torsor-moduli type $\BG\CHunits$, its canonical point, and the underlying
type equivalence of Equation~\eqref{eq:main-result}.  The generic theory also
establishes connectedness and one-truncation for the torsor moduli, though the
constructive-complex module uses only the basepoint and loop equivalence.

\subsection{Explicit non-results}
The following statements are outside the current formal endpoint.
\begin{enumerate}[label=\arabic*.]
\item No coordinate-level apartness/nonzero characterization of algebraic
complex units is proved.
\item No group-level loop recovery theorem is packaged for
\agda{complexUnitsLoopEquiv}.
\item No general proof identifies complex-unit multiplication with loop
concatenation for that displayed equivalence.
\item No dependent-family transport computation or monodromy calculation is
constructed.
\item No quantum-gate realization is implemented.
\item Myers--Sati--Schreiber Definition~6.7 is not formalized.
\end{enumerate}

These are not cosmetic omissions.  A group isomorphism requires a stated loop
group structure and proofs preserving product, unit, and inverse.  A quantum
transformation requires a specified dependent family, a loop or higher action
on it, and an analysis of the corresponding transport.  Neither structure
follows merely from the type equivalence already proved.

\subsection{Next mathematical stages}
A later development may choose a constructive apartness relation for complex
numbers and prove its relation to algebraic units.  It may then formulate a
group-level refinement of loop recovery with explicit orientation conventions.
Only after specifying a dependent family over a candidate delooping could the
project study transport, its computational behaviour, and its possible
relationship with the topological-quantum motivation of \cite{mss}.  Those
would be new formalization tasks rather than consequences of this paper.

\begin{thebibliography}{9}
\bibitem{hottbook}
The Univalent Foundations Program.
\newblock \emph{Homotopy Type Theory: Univalent Foundations of Mathematics}.
\newblock Institute for Advanced Study, 2013.
\newblock \url{https://homotopytypetheory.org/book/}.

\bibitem{brough}
Jackson Brough.
\newblock \emph{hott-reals}, thesis branch, revision
\texttt{9fcc142b86d0c1e2bd6d06f86f7617200ad74857}.
\newblock \url{https://github.com/utahplt/hott-reals}.

\bibitem{cubical}
Cubical Agda contributors.
\newblock \emph{Cubical Agda standard library}, v0.9-compatible line.
\newblock Pinned project checkout:
\texttt{87af036ebd1c87956cf08fea44e62c7cf76392ca}.
\newblock \url{https://github.com/agda/cubical}.

\bibitem{agda}
The Agda Team.
\newblock \emph{Agda 2.8.0-3d04bac}.
\newblock \url{https://agda.readthedocs.io/}.

\bibitem{mss}
David Jaz Myers, Hisham Sati, and Urs Schreiber.
\newblock \emph{Topological Quantum Gates in Homotopy Type Theory}.
\newblock \emph{Communications in Mathematical Physics} \textbf{405},
Article 172 (2024).
\newblock \url{https://doi.org/10.1007/s00220-024-05020-8}.
\end{thebibliography}
